% Unregistered D.65 candidate. The intrinsic five-adic topology below
% is deliberately distinguished from continuous-cohomology topology.
\subsection{The abstract module carried by the nonseparated residue kernel}

\begin{statement}{D.65}{Proved}{The continuum-sized complete module in every finite-ramification Jannsen kernel}
Let \(s\) be any finitely ramified Alexander twist, and
retain \(S\), \({\cal C}_S\), \(D_{s,q}\), and the compact
group \({\cal L}_s\) of D.48. Put
\[
 \begin{split}
 {\cal R}_s&=
  \frac{{\cal C}_S\cap\prod_{q\notin S}D_{s,q}A_s}
       {D_s{\cal C}_S},\\
 {\cal J}_s&=\varprojlim\nolimits_r^1
                         H^1(G_{\mathbb Q},A_s/\mathfrak m^rA_s).
 \end{split}                                             \tag{D.65.1}
\]
For a set \(I\), define the complete free module
\[
 \widehat{\bigoplus}_{i\in I}\mathbb Z_5
   :=\varprojlim_n\bigoplus_{i\in I}\mathbb Z/5^n
   =\left\{(a_i)\in\prod_I\mathbb Z_5:
       \forall n,\ a_i\in5^n\mathbb Z_5
                         \text{ for all but finitely many }i\right\}.
                                                        \tag{D.65.2}
\]
There are noncanonical isomorphisms of abstract
\(\mathbb Z_5\)-modules
\[
 \begin{split}
 {\cal R}_s&\simeq
              \widehat{\bigoplus}_{2^{\aleph_0}}\mathbb Z_5,\\
 {\cal J}_s&\simeq {\cal L}_s\oplus
              \widehat{\bigoplus}_{2^{\aleph_0}}\mathbb Z_5.
 \end{split}                                             \tag{D.65.3}
\]
In particular the residue kernel is torsion-free, is
complete and separated for its intrinsic five-adic
topology, and has
\[
           \dim_{\mathbb F_5}({\cal R}_s/5{\cal R}_s)
                                      =2^{\aleph_0}.
                                                        \tag{D.65.4}
\]
The first isomorphism in (D.65.3) is also a homeomorphism
for these intrinsic five-adic topologies.
Equivalently its right-hand complete free module is
abstractly isomorphic to \(\prod_{\mathbb N}\mathbb Z_5\).
Under this equivalence the intrinsic topology on the
product is the uniform five-adic topology, not its
usual compact product topology.

These are not the subspace topologies from continuous
cohomology. In that topology \({\cal J}_s\) is the
closure of zero, as in D.62. The second splitting in
(D.65.3) is only an abstract module splitting; it does
not replace any cohomology group by a separated
quotient or assert a continuous retraction onto the
compact group with its original topology.

For the families of D.44 and D.64,
\({\cal L}_s=0\), so (D.65.3) gives the complete
abstract isomorphism type of their Jannsen kernel.
The same complete module is the kernel in each of
the three Tate weights of D.1. For a general finitely
ramified twist, the compact term \({\cal L}_s\)
is retained. This result makes no classification
claim for the whole kernel of an infinitely
ramified twist.
\end{statement}

\begin{proof}
We give the module and completeness arguments because
the two topologies here have different purposes.
Identify \(A_s\) additively with
\(\Lambda=\mathbb Z_5[[u,v]]\), and write
\[
 \begin{split}
 P&=\prod_{q\notin S}\Lambda,\qquad C={\cal C}_S,\\
 H&=\{(a_q)\in P:(D_{s,q}a_q)_q\in C\}.
 \end{split}                                             \tag{D.65.5}
\]
Each \(D_{s,q}\) is injective by D.48. Taking its unique
coordinatewise inverse on the image therefore identifies
\[
                             {\cal R}_s=H/C             \tag{D.65.6}
\]
as abstract modules. Also \(C\subset H\): every operator
\(D_{s,q}\) preserves \(\mathfrak m^n\), uniformly in
\(q\), because the actual arithmetic action does.

\paragraph{A stronger complete topology for the module calculation.}
For this paragraph only, give \(P\) the uniform
five-adic topology, whose neighborhoods are \(5^nP\).
The ring \(\Lambda\) is complete and separated for its
five-adic topology, so \(P\) is complete for this
uniform topology. The module \(C\) is closed in it.
Indeed if \(a\) can be approximated by elements of
\(C\) modulo every \(5^nP\), fix \(n\) and choose one
such approximation \(c\). Away from a finite set,
\(c_q\in\mathfrak m^n\), and
\(a_q-c_q\in5^n\Lambda\subset\mathfrak m^n\).
Thus \(a\in C\).

The coordinatewise operator \(D_s:P\to P\) is
continuous for the uniform five-adic topology:
each coordinate is \(\mathbb Z_5\)-linear and preserves
\(5^n\Lambda\), with the same bound for all coordinates.
Hence \(H=D_s^{-1}(C)\) is closed and complete as well.
Both \(C\) and \(H\) are five-saturated in \(P\).
For \(C\), use the exact order formula
\[
 \operatorname{ord}_{\mathfrak m}(5^n z)
          =n+\operatorname{ord}_{\mathfrak m}(z)
                        \quad(z\ne0).                  \tag{D.65.7}
\]
It follows either from coefficients or from the
polynomial integral-domain associated graded ring
of \(\mathbb Z_5[[u,v]]\) at \((5,u,v)\).
Thus \((5^na_q)_q\) tends to zero in \(\mathfrak m\)
if and only if \((a_q)_q\) does. For \(H\), apply the
same observation to \(D_s(5^na)=5^nD_s(a)\).
Consequently
\[
 H\cap5^nP=5^nH,\qquad C\cap5^nP=5^nC,
                  \qquad H/C\text{ is torsion-free}.
                                                        \tag{D.65.8}
\]
The topology on \(H\) induced from \(P\) is therefore
its intrinsic five-adic topology, and its quotient
topology on \(H/C\) has neighborhoods precisely
\(5^n(H/C)\).

This quotient is complete and separated. Separation
is the closedness of \(C\) just proved. For completeness,
take a compatible system of classes modulo
\(5^n(H/C)\). Choose representatives \(h_n\in H\)
successively. Compatibility lets one change the next
representative by an element of \(C\) so that
\(h_{n+1}-h_n\in5^nH\). Their limit exists in the
complete module \(H\), and its class gives the required
limit. This proves the asserted completeness directly;
it does not presume completeness of a quotient for
the original cohomology topology.

\paragraph{The exact number of five-adic generators.}
The Chebotarev construction in the last paragraph
of D.48 supplies distinct primes \(q_j\notin S\)
with
\[
                         D_{s,q_j}(1)\in\mathfrak m^j.
                                                        \tag{D.65.9}
\]
For any \(a=(a_j)\in\prod_j\mathbb Z_5\), put
\(h(a)_{q_j}=a_j\), with all other coordinates zero.
Then \(h(a)\in H\). Its class in \(H/C\) is the
copy of \(\prod_j\mathbb Z_5/c_0(\mathbb N,\mathbb Z_5)\)
already constructed there. It remains injective
after reduction modulo five. Indeed if
\(h(a)=5b+c\), with \(b\in H\), \(c\in C\), evaluation
at \(u=v=0\) gives
\[
                 a_j=5b_{q_j}(0,0)+c_{q_j}(0,0).
\]
The last constants tend to zero five-adically,
so \(a_j\bmod5\) is zero for all sufficiently large
\(j\). This is exactly the condition that \(a\)
be zero in the reduction modulo five of that
product-by-\(c_0\) quotient. Thus
\[
 \mathbb F_5^{\mathbb N}/\mathbb F_5^{(\mathbb N)}
             \lhook\joinrel\longrightarrow (H/C)/5(H/C).
                                                        \tag{D.65.10}
\]
The vector space on the left has dimension
\(2^{\aleph_0}\): over a finite field an infinite
dimensional vector space has cardinality equal
to its dimension, and quotienting by a countable
dimensional subspace does not change this dimension.
On the other hand \(P\), hence \(H/C\), has cardinality
at most \(2^{\aleph_0}\). This proves (D.65.4).

For completeness, a separated complete torsion-free
\(\mathbb Z_5\)-module \(M\) is the complete free module
on a basis of \(M/5M\). Choose arbitrary lifts of that
basis. A family of coefficients tending to zero
five-adically has a well-defined sum in \(M\).
Successively lifting the finitely supported expansion
of each residual class modulo five proves that these
sums exhaust \(M\). To check injectivity, take the
smallest valuation of a nonzero coefficient family,
divide by that power of five, and reduce modulo five;
the resulting nonempty finite basis combination is
nonzero. The map is an isomorphism and preserves
every subgroup \(5^nM\), hence is a homeomorphism
for the intrinsic topologies. Apply this to
\(M=H/C\), using (D.65.4), to obtain the first line
of (D.65.3).
The module \(\prod_{\mathbb N}\mathbb Z_5\) is also
intrinsically complete, separated, and torsion-free,
and its reduction is \(\mathbb F_5^{\mathbb N}\),
of the same continuum dimension. The same classification
proves the equivalent countable-product description.

\paragraph{Splitting off the retained compact term.}
D.48 gives the actual exact sequence of modules
\[
       0\longrightarrow{\cal L}_s
         \longrightarrow{\cal J}_s
         \longrightarrow{\cal R}_s\longrightarrow0.
                                                        \tag{D.65.11}
\]
Since the last module is torsion-free, the first
inclusion is pure: for every \(n\),
\({\cal L}_s\cap5^n{\cal J}_s=5^n{\cal L}_s\).
We spell out why compactness of \({\cal L}_s\), with
its original compact-module topology, gives an
abstract retraction. Introduce a variable \(x_j\)
with values in \({\cal L}_s\) for each
\(j\in{\cal J}_s\), and impose all equations
\[
 x_{j+j'}=x_j+x_{j'},\qquad
 x_{\lambda j}=\lambda x_j\ (\lambda\in\mathbb Z_5),
                 \qquad x_l=l\ (l\in{\cal L}_s).
                                                        \tag{D.65.12}
\]
Each finite subsystem, viewed as equations with
constants in \({\cal L}_s\), has the evident solution
\(x_j=j\) in \({\cal J}_s\). Smith normal form over
the discrete valuation ring \(\mathbb Z_5\) reduces
solvability of any such finite linear system to
divisibility of constants by powers of five and
zero equations. Purity makes those divisibilities
solvable already in \({\cal L}_s\). Thus every finite
subsystem has a solution in that compact module.
All equations are closed in the compact product
\({\cal L}_s^{{\cal J}_s}\). The finite-intersection
property gives a simultaneous solution, hence an
abstract \(\mathbb Z_5\)-linear retraction of the
inclusion in (D.65.11). This proves its splitting
and the second line of (D.65.3). No continuity of
this retraction for the nonseparated cohomology
topology has been asserted.

For D.44 the compact finite-support second cohomology
is zero, and for D.64 its maps to unrestricted finite
cohomology are injective. In both cases D.48's
\({\cal L}_s\) is therefore zero. Finally the three
Tate kernels in D.1 have the same proof with
\(\Lambda=\mathbb Z_5\), \(\mathfrak m=(5)\), and
\(D_q=q^{j-1}-1\). These operators are injective,
their arbitrarily divisible unit-input values are
given by the primes used in D.1, and there is no
compact base term. Thus their complete abstract
module type is the same one in (D.65.3).
\end{proof}
